
\documentclass[11pt]{article}

\usepackage{amsmath}
\usepackage{amssymb}
\usepackage{graphicx}
\usepackage{hyperref}
\usepackage{geometry}
\geometry{margin=1in}

\title{The Prime Exponent Tree (PET): \\
A Structural View of Integer Factorization}

\author{Giancarlo Comneno \\
\small Independent Research}

\date{\today}

\begin{document}

\maketitle

\begin{abstract}
We introduce the \textbf{Prime Exponent Tree (PET)}, a recursive tree representation
of integers derived from prime factorization. In this model, every integer is
represented as a tree whose nodes correspond to prime bases and whose children
recursively encode the structure of the corresponding exponents.

Using a complete dataset of all integers up to $10^6$, we empirically explore the
space of PET structures. Surprisingly, despite the large number of integers,
only a small number of distinct structural shapes appear. Among the first one
million integers, we observe only 78 distinct structural shapes, and the
distribution of these shapes is highly concentrated: seven shapes account for
approximately 87\% of all integers.

We further measure the structural entropy of the PET distribution and observe
that the complexity of integer multiplicative structure grows approximately as
$\log\log N$. These results suggest that the multiplicative structure of integers
is dramatically simpler than their numerical magnitude might suggest.
\end{abstract}

\section{Introduction}

Integers grow rapidly, but their multiplicative structure evolves much more
slowly. Classical results in analytic number theory such as the
Hardy--Ramanujan theorem and the Erd\H{o}s--Kac theorem show that the number
of distinct prime factors of an integer grows on average like $\log\log n$.

However, traditional approaches to integer factorization focus primarily on the
values of primes rather than on the \emph{structure} of factorizations.
In this work we explore a complementary viewpoint: representing integers as
recursive combinatorial objects.

We introduce the \textbf{Prime Exponent Tree (PET)}, a canonical representation
of integers in which prime bases appear as nodes and the structure of exponents
is recursively encoded as subtrees.

For example:

\begin{equation}
12 = 2^2 \cdot 3
\end{equation}

is represented as a tree where the node corresponding to the prime $2$
contains a subtree representing the integer $2$.

This representation allows us to study the \emph{structural space} of
integer factorizations independently of the actual prime values.

Our central empirical question is therefore:

\begin{quote}
How many distinct multiplicative \emph{structures} appear among the integers
up to $N$?
\end{quote}

Using a complete dataset of all integers up to one million, we show that the
space of such structures is surprisingly small and highly organized.

The remainder of this paper explores this structural space using the PET
representation and presents several empirical observations about its growth,
distribution, and complexity.

\section{The Prime Exponent Tree}

(Section to be expanded in the next draft.)

\section{Dataset and Experimental Setup}

(To be completed.)

\section{Structural Shape Distribution}

(To be completed.)

\section{Entropy of Integer Structure}

(To be completed.)

\section{Discussion}

(To be completed.)

\section{Conclusion}

(To be completed.)

\end{document}
